\section{Related work}
\label{sec:related}

\paragraph{Personal identity and its thought experiments.}
The question whether a copy, a teleported replica or a gradually replaced
system is the same subject has a long literature
\citep{Williams1970, Parfit1971, Perry1972, Lewis1976, Nozick1981, Parfit1984}.
\citet{Olson2023SEP} surveys the fission problem and the two standard replies.
The reduplication argument, that a relation which could hold between one
earlier and two later persons cannot by itself make them identical, goes back
to \citet{Williams1957}.
On multiple occupancy, two persons share their stages before the division
\citep{Lewis1976}. On non-branching criteria, a person survives only if no
equally good competitor exists. Parfit discusses criteria of this kind, and
Nozick's closest-continuer theory is a distinct proposal under which a tie
leaves no continuer \citep{Nozick1981}. We do not take a position among these
views. Our structural results sit beside them: under fission the continuation
relation has two values, and no rule respecting the symmetry selects one
(Section~\ref{sec:branching}). Perry's notion of a branch, a maximal set of
mutually related stages \citep{Perry1972}, invites comparison with our strict
lineage. The two differ: our chains are checked sequences of handoffs with no
maximality requirement. \citet{Gustafsson2019} asks what form a non-branching
clause should take when persistence rests on a continuity relation that can
hold from one to many, and treats fusion as well as fission. Our relation
holds between finite systems, is checked by a proof assistant, and is not
analysed as identity.

\paragraph{Uploading and gradual replacement.}
\citet{Chalmers2010Singularity} separates the question of organization from
the question of identity, and distinguishes gradual, destructive and
reconstructive uploading. \citet{Chalmers1995Qualia} uses gradual replacement
to argue for organizational invariance of experience.
\citet{Cerullo2015} argues informally that uploading can yield two
continuers with no reason to favour either, under a criterion of branching
identity. \citet{Bauer2017} argues against branching identity.
\citet{CorabiSchneider2012} ask whether an upload is at least a continuation of
the original, in terms of its spatiotemporal and causal connection to it. In
reply, \citet{Chalmers2012Reply} notes, with the example of a university that
moved, that a moved thing may for a time exist at both places or at neither.
He also notes that a gradually uploaded system stays connected to the
original throughout. \citet{Weber2025Selective} argues that gradual replacement
and scan-and-copy differ metaphysically, in the immanent causation they
involve. Our three questions formalize
questions of this kind on finite systems, and we too keep identity out of
them. We do not claim that gradual versus instant uploading corresponds to our
distinction between overlap and direct succession. Our classes require
operational checks that timing alone does not supply. And nothing here bears
on whether an upload is conscious.

\paragraph{Temporal experience.}
Overlap models of the stream of consciousness, in which successive specious
presents \citep[ch.~XV]{James1890} share content, are the phenomenological
counterpart of our live-overlap handoffs \citep{Dainton2000, Dainton2023SEP}.
\citet{DaintonBayne2005} propose that phenomenal continuity, rather than
psychological continuity, is the guide to personal persistence, and
\citet{Dainton2008} develops that view and treats fission and fusion within
it. Our overlaps are shared system
state, not co-consciousness, and we adopt no model of experienced duration. We
also take no side between views on which experience inherits the temporal
structure of its objects \citep{Phillips2014} and views on which it need not
\citep{Lee2014}. Within integrated information theory,
\citet{ComolattiGrassoTononi2025} account for experienced time by moments that
overlap in a directed way, and take the present to be a cause--effect
structure specified by a system in its current state. Our predictive present
is a different object, a field of the temporal signature carried across a
transfer. We establish no bridge between the two and draw no conclusion about
their account.

\paragraph{Mathematical consciousness science.}
Our setting is that of mathematical models of consciousness, where a model's
objects and their well-definedness conditions are stated explicitly
\citep{Kleiner2020MathematicalModels}. In that framework experience is
defined at an instant, and what counts as an instant is left to the model
\citep[Definition~2]{Kleiner2020MathematicalModels}. Our checks concern a
different object: how a temporal signature, read off each state, is carried
and maintained over handoff episodes and future sequences of actions. Integrated information theory
requires a unique maximal substrate. When the maximal candidates tie and all
of them overlap, its stated rule picks the next-best unique system
\citep{AlbantakisEtAl2023IIT4}. That the theory's maximizations need not have
unique results was already noted in its formalization
\citep[Section~12.2]{KleinerTull2021}. This is a
case from the field's own literature of a theory facing a uniqueness
requirement over tied candidates. It is a different move from selecting among
the tied ones, and our results do not show that the rule needs to break any
symmetry. Substrate questions also appear as formal no-go arguments
\citep{KleinerLudwig2024Neurons}. Those concern which systems are conscious,
whereas ours concern which systems count as structural continuers, so the two
neither support nor contradict each other.

\paragraph{Mathematical precursors.}
The ingredients are standard, and we credit them where they are used:
\begin{itemize}
\item gluing over covers with running intersection (acyclic database schemes,
\citealp[Theorem~3.4]{BeeriEtAl1983}; see also \citealp{Vorobev1962} and
\citealp{Abramsky2013}). Global sections and their obstruction are treated in
\citet{AbramskyBrandenburger2011}, whose obstructions concern distributions. Our
cyclic obstruction is instead a failure of running intersection.
\item confluence and normal forms \citep{Newman1942, BaaderNipkow1998}.
\item bisimulation as behavioural equivalence \citep{Park1981, Milner1989},
and its probabilistic form \citep{LarsenSkou1991}.
\item lumpability of Markov chains \citep{KemenySnell1976}.
\item controlled invariance and viability \citep{Blanchini1999, Aubin1991},
and traps in infinite games \citep{Mazala2002}.
\item safety and liveness properties \citep{AlpernSchneider1985}.
\item limits of finite-step reachability \citep{Bertsekas1972}.
\item holonomy \citep{Nakahara2003}.
\item predictive equivalence \citep{ShaliziCrutchfield2001} and predictive
representations of state \citep{LittmanSuttonSingh2002}.
\end{itemize}
We claim no priority for any of these ideas, separately or together.

\paragraph{The Six Birds framework.}
The questions and the design of the models come from a broader programme, the
Six Birds framework, which we describe in Appendix~\ref{app:sixbirds} together
with the sources the design drew on. None of our results depends on it.
